% ============================================================================
%  CORRECTED 2026-08-17.  Two edits, both required before this exports to a
%  PDF that can be served from ecowoods.ca:
%
%  1. TITLE STRIP.  Was "Est. 1998 / 2000 | 5,200+ Homes".  Both strings are
%     on the retired-claims list in scripts/verify-business-facts.mjs — the
%     home count is invented and the founding year is BUSINESS_NAP.foundedYear
%     (2000).  They were removed from the whole codebase in the business-facts
%     remediation, and the guard now reads PDF text as well, so a PDF carrying
%     them turns the build red on purpose.
%
%  2. FOOTER.  Was "Vincenzo Ceccarelli Grimaldi | Grid Networks Engineer |
%     LinkedIn | GitHub" on every page.  A homeowner reading about belt sanders
%     does not need a Grid Networks Engineer's GitHub, and it quietly
%     reassigns authorship of an EcoWoods document to a different entity in a
%     different industry.
%
%  Re-export from Overleaf with the four machine .jpg files in the same folder,
%  then:  mv docs/papers-pending/*.pdf apps/web/public/papers/
%  The download buttons appear automatically — pdfIsPublished() checks the file.
% ============================================================================

\documentclass[aspectratio=169,10pt]{beamer}
\usepackage{textcomp}
\usetheme{Madrid}
\usecolortheme{default}
\usefonttheme{professionalfonts}
\usepackage[utf8]{inputenc}
\usepackage[T1]{fontenc}
\usepackage{booktabs}
\usepackage{tikz}
\usetikzlibrary{shapes.geometric,arrows.meta,positioning,calc,fit,backgrounds,shadows}
\usepackage{graphicx}
\usepackage{hyperref}
\usepackage{amsmath}
\usepackage{xcolor}
\usepackage{textcomp}

% ========== ECOWOODS COLOR SYSTEM ==========
\definecolor{ecodark}{RGB}{16,20,28}
\definecolor{ecowood}{RGB}{145,95,45}
\definecolor{ecoaccent}{RGB}{0,100,120}
\definecolor{ecolite}{RGB}{250,247,242}
\definecolor{ecolightblue}{RGB}{220,232,242}
\definecolor{ecogray}{RGB}{65,75,85}

\setbeamercolor{structure}{fg=ecoaccent}
\setbeamercolor{frametitle}{bg=ecodark,fg=white}
\setbeamercolor{title}{fg=ecodark}
\setbeamercolor{block title}{bg=ecodark,fg=white}
\setbeamercolor{block body}{bg=ecolite}
\setbeamercolor{alerted text}{fg=ecowood}
\setbeamercolor{block title alerted}{bg=ecowood,fg=white}
\setbeamercolor{block body alerted}{bg=ecolite}
\setbeamercolor{author in head/foot}{bg=ecolightblue,fg=ecodark}

\setbeamertemplate{navigation symbols}{}
\setbeamertemplate{footline}{%
  \leavevmode%
  \hbox{%
  \begin{beamercolorbox}[wd=\paperwidth,ht=5ex,dp=3.2ex,leftskip=1.2em,rightskip=1.2em]{author in head/foot}%
    \scriptsize\color{ecodark}
    \textbf{EcoWoods} \quad|\quad Hardwood, Done Once. Done Right. \quad|\quad
    \href{https://ecowoods.ca}{ecowoods.ca} \quad|\quad (647) 244-5156
    \hfill
    \insertframenumber/\inserttotalframenumber
  \end{beamercolorbox}%
  }%
  \vskip0pt%
}

\newcommand{\framedimage}[2][]{%
  \begin{tikzpicture}
    \node[inner sep=1.2pt, draw=ecodark!30, line width=0.7pt, rounded corners=1.8pt] {
      \includegraphics[#1]{#2}
    };
  \end{tikzpicture}%
}

\begin{document}

% ========== 1. TITLE ==========
\begin{frame}
\begin{center}
\vspace{0.3cm}
{\fontsize{22}{26}\selectfont\textbf{\color{ecodark}THE CRAFT OF HARDWOOD FLOORING}}\\[0.12cm]
{\fontsize{17}{21}\selectfont\textbf{\color{ecodark}Machines, Process \& Mastery}}\\[0.35cm]
{\large An Educational Framework for Understanding\\How Professional Hardwood Companies Actually Work}\\[0.4cm]
{\large\textit{Hardwood, Done Once. Done Right.}}\\[0.22cm]
{\small ecowoods.ca \quad|\quad Est. 2000 \quad|\quad Toronto \& the GTA \quad|\quad Lifetime Workmanship Warranty}
\end{center}
\end{frame}

% ========== 2. LEARNING OBJECTIVES ==========
\begin{frame}{Learning Objectives}
\vspace{0.05cm}
\begin{enumerate}
  \item What a professional hardwood flooring company actually does
  \item The EcoWoods difference: protocol over speed
  \item The four core machines that define the craft
  \item Machine 1 — Belt Floor Sander (the big machine)
  \item Machine 2 — Floor Edger
  \item Machine 3 — Planetary / Multi-Disc Sander
  \item Machine 4 — Floor Buffer \& Screening
  \item How the full sequence produces a floor that lasts decades
  \item Why equipment + people + process = mastery
\end{enumerate}
\end{frame}

% ========== 3. WHAT COMPANIES DO ==========
\begin{frame}{What Hardwood Flooring Companies Actually Do}
\begin{columns}[T]
\begin{column}{0.52\textwidth}
\begin{itemize}
  \item Assess the house (moisture, substrate, RH history)
  \item Specify the correct product (solid vs engineered)
  \item Install or fully refinish with progressive sanding
  \item Apply finish systems that perform for 25–35+ years
  \item Protect the client from change orders and callbacks
\end{itemize}
\end{column}
\begin{column}{0.45\textwidth}
\begin{alertblock}{Core Truth}
A hardwood floor is a living system.\\
The company that masters both the forest and the machines becomes the only rational choice.
\end{alertblock}
\end{column}
\end{columns}
\vspace{0.25cm}
\begin{center}
Most market players optimise for lowest bid and speed.\\
Master companies optimise for zero callbacks and lifetime performance.
\end{center}
\end{frame}

% ========== 4. ECOWOODS DIFFERENCE ==========
\begin{frame}{The EcoWoods Standard}
\begin{itemize}
  \item Fixed price in writing — no open-ended change orders
  \item Zero dust (99.7\% HEPA containment — Festool \& Bona systems)
  \item Only salaried master artisans (many 10+ years) — never subcontractors
  \item Moisture testing + documented protocol on every job
  \item Lifetime workmanship warranty in the contract
  \item Ability and willingness to refuse jobs that will fail
\end{itemize}
\vspace{0.2cm}
\begin{block}{Educational Principle}
Restraint is a competitive advantage.\\
The machines are only as good as the people and the protocol that control them.
\end{block}
\end{frame}

% ========== 5. THE FOUR MACHINES OVERVIEW ==========
\begin{frame}{The Four Core Machines}
\begin{center}
\begin{tabular}{cl}
\toprule
\textbf{\#} & \textbf{Machine} \\
\midrule
1 & Belt Floor Sander (Big Machine) — field leveling \& finish removal \\
2 & Floor Edger — walls, corners, stairs, tight spaces \\
3 & Planetary / Multi-Disc Sander — refining \& blending \\
4 & Floor Buffer / Screening Machine — final prep \& intercoat abrasion \\
\bottomrule
\end{tabular}
\end{center}
\vspace{0.3cm}
\begin{center}
These four machines, used in sequence with progressive grits and proper dust control,\\
are the mechanical backbone of every professional hardwood refinishing or site-finished installation.
\end{center}
\end{frame}

% ========== 6. MACHINE 1 ==========
\begin{frame}{Machine 1 — Belt Floor Sander}
\begin{columns}[T]
\begin{column}{0.48\textwidth}
\begin{itemize}
  \item Primary ``big machine''
  \item Continuous abrasive belt over a drum ($\approx$200\,mm)
  \item Removes old finish, levels the floor
  \item Progressive grits: 36 $\rightarrow$ 60 $\rightarrow$ 80/100
  \item Walk-behind, high material removal rate
\end{itemize}
\vspace{0.12cm}
\begin{alertblock}{Key Skill}
Keep it moving. Stopping creates divots. Pressure and tracking determine quality.
\end{alertblock}
\end{column}
\begin{column}{0.48\textwidth}
\begin{center}
\framedimage[width=0.95\textwidth]{01_belt_sander.jpg}
\end{center}
\end{column}
\end{columns}
\end{frame}

% ========== 7. MACHINE 1 DETAIL ==========
\begin{frame}{Belt Sander — Why It Matters}
\begin{itemize}
  \item Handles $\approx$80\% of material removal on a typical refinish
  \item Creates the flat plane that every subsequent machine builds upon
  \item Incorrect use leaves chatter, waves, or side-cut marks that are extremely difficult to erase
  \item At EcoWoods: always paired with HEPA extraction; operated only by salaried artisans
\end{itemize}
\vspace{0.25cm}
\begin{block}{Educational Takeaway}
The belt sander is the foundation of the mechanical process.\\
Mastery here determines whether the floor will look and perform correctly for decades.
\end{block}
\end{frame}

% ========== 8. MACHINE 2 ==========
\begin{frame}{Machine 2 — Floor Edger}
\begin{columns}[T]
\begin{column}{0.48\textwidth}
\begin{itemize}
  \item High-speed rotating disc
  \item Reaches every place the belt cannot
  \item Walls, baseboards, closets, stairs, under cabinets
  \item Must match the grit sequence of the field
  \item Most aggressive of the four machines
\end{itemize}
\vspace{0.12cm}
\begin{alertblock}{Key Skill}
Technique is everything. Swirl marks from poor edging are highly visible under finish.
\end{alertblock}
\end{column}
\begin{column}{0.48\textwidth}
\begin{center}
\framedimage[width=0.95\textwidth]{02_edger.jpg}
\end{center}
\end{column}
\end{columns}
\end{frame}

% ========== 9. MACHINE 2 DETAIL ==========
\begin{frame}{Edger — Why It Matters}
\begin{itemize}
  \item Edges are where most low-quality jobs fail visually
  \item Must be blended carefully into the field (never edge too far out)
  \item Corners and detail areas often still require hand work after the edger
  \item Dust control is harder than on the big machine — HEPA connection is non-negotiable
\end{itemize}
\vspace{0.25cm}
\begin{block}{Educational Takeaway}
The edger reveals the true skill level of the crew.\\
At EcoWoods, salaried artisans treat the perimeter with the same care as the open field.
\end{block}
\end{frame}

% ========== 10. MACHINE 3 ==========
\begin{frame}{Machine 3 — Planetary / Multi-Disc Sander}
\begin{columns}[T]
\begin{column}{0.48\textwidth}
\begin{itemize}
  \item Multiple counter-rotating discs (Trio-style)
  \item Creates random, multi-directional scratch pattern
  \item Erases belt lines and edger swirls
  \item Excellent for flattening and blending
  \item Forgiving yet highly effective
\end{itemize}
\vspace{0.12cm}
\begin{alertblock}{Key Skill}
This is the refining stage. It turns an aggressive cut into a uniform surface ready for finish.
\end{alertblock}
\end{column}
\begin{column}{0.48\textwidth}
\begin{center}
\framedimage[width=0.95\textwidth]{03_planetary_sander.jpg}
\end{center}
\end{column}
\end{columns}
\end{frame}

% ========== 11. MACHINE 3 DETAIL ==========
\begin{frame}{Planetary Sander — Why It Matters}
\begin{itemize}
  \item Separates a ``good'' sanding job from a master-level result
  \item Ideal for multi-species floors and soft/hard grain differences
  \item Prepares the surface so stain and finish absorb evenly
  \item Modern dust-free systems integrate planetary action with high-efficiency extraction
\end{itemize}
\vspace{0.25cm}
\begin{block}{Educational Takeaway}
Skipping or rushing the planetary pass is one of the most common reasons floors show machine marks after finishing.
\end{block}
\end{frame}

% ========== 12. MACHINE 4 ==========
\begin{frame}{Machine 4 — Floor Buffer \& Screening}
\begin{columns}[T]
\begin{column}{0.48\textwidth}
\begin{itemize}
  \item Large circular drive plate (16–20'')
  \item Screens the floor with fine mesh (100–150 grit)
  \item Creates uniform microscopic scratch for adhesion
  \item Used between coats (intercoat abrasion)
  \item Final mechanical step before finish
\end{itemize}
\vspace{0.12cm}
\begin{alertblock}{Key Skill}
Screening is light abrasion only — not aggressive sanding. Technique prevents burnishing or residual swirl.
\end{alertblock}
\end{column}
\begin{column}{0.48\textwidth}
\begin{center}
\framedimage[width=0.95\textwidth]{04_buffer.jpg}
\end{center}
\end{column}
\end{columns}
\end{frame}

% ========== 13. MACHINE 4 DETAIL ==========
\begin{frame}{Buffer — Why It Matters}
\begin{itemize}
  \item Last chance to create a perfect substrate for the finish system
  \item Mandatory between coats on multi-coat water-based finishes
  \item At EcoWoods: always performed inside the full HEPA dust-containment protocol
  \item Enables the low-VOC, GreenGuard Gold finishes we specify
\end{itemize}
\vspace{0.25cm}
\begin{block}{Educational Takeaway}
The buffer does not correct bad sanding — it prepares a correctly sanded floor for decades of performance.
\end{block}
\end{frame}

% ========== 14. FULL SEQUENCE ==========
\begin{frame}{The Full Professional Sequence}
\begin{enumerate}
  \item Moisture testing \& acclimation (minimum 72 hours in the actual space)
  \item Belt sander — progressive grits, field only
  \item Edger — matching grits on all perimeters and details
  \item Planetary / multi-disc — refine and blend field + edges
  \item Buffer / screening — final uniform surface
  \item Vacuum, tack, apply finish system (water-based, low-VOC)
  \item Intercoat screening with buffer between coats
  \item Final coat + documented lifetime workmanship warranty
\end{enumerate}
\vspace{0.15cm}
\begin{center}
\textbf{Any skipped step is a future liability.}
\end{center}
\end{frame}

% ========== 15. WHY THIS MATTERS ==========
\begin{frame}{Equipment + People + Process = Mastery}
\begin{columns}[T]
\begin{column}{0.48\textwidth}
\begin{block}{Market Average}
\begin{itemize}
  \item Lowest bid
  \item Subcontractors
  \item Optional moisture testing
  \item Speed over protocol
\end{itemize}
\end{block}
\end{column}
\begin{column}{0.48\textwidth}
\begin{alertblock}{EcoWoods}
\begin{itemize}
  \item Fixed price
  \item Salaried master artisans
  \item Moisture testing as first gate
  \item Zero-callback protocol
\end{itemize}
\end{alertblock}
\end{column}
\end{columns}
\vspace{0.25cm}
\begin{center}
The machines are available to everyone.\\
The combination of machines + trained people + documented restraint is rare.
\end{center}
\end{frame}

% ========== 16. ACTION FOR STUDENTS ==========
\begin{frame}{How to Use This Material}
\begin{enumerate}
  \item Study each machine image and explanation from the accompanying zip
  \item Film short educational clips: one machine per video, showing purpose and correct use
  \item Emphasize that the machines only produce mastery when controlled by protocol and skilled artisans
  \item Link every demonstration back to the EcoWoods standard: fixed price, zero dust, lifetime warranty
\end{enumerate}
\vspace{0.25cm}
\begin{block}{Final Teaching Point}
The floor will declare the competence of the people and the process that created it.\\
Make sure it only ever declares mastery.
\end{block}
\end{frame}

% ========== 17. KEY TAKEAWAYS ==========
\begin{frame}{Key Educational Takeaways}
\begin{enumerate}
  \item Four machines form the mechanical core of professional hardwood work
  \item Sequence and progressive grits are non-negotiable
  \item Dust control (HEPA) is part of modern professional standard
  \item People and protocol matter more than any single machine
  \item EcoWoods integrates all of the above into a fixed-price, lifetime-warrantied system
\end{enumerate}
\end{frame}

% ========== 18. FINAL ==========
\begin{frame}{Final Mandate}
\begin{center}
\Huge
The floor will declare your competence.\\[0.35em]
{\large Make sure it only ever declares mastery.}
\end{center}
\vspace{0.5cm}
\begin{center}
{\LARGE\textbf{EcoWoods}}\\[0.1em]
Hardwood, Done Once. Done Right.\\[0.22em]
{\small ecowoods.ca \quad · \quad (647) 244-5156}\\[0.08em]
{\footnotesize 32 Norfield Crescent, Toronto}
\end{center}
\end{frame}

\end{document}
